\subsection{Thiết kế pipeline hỏi đáp pháp lý bằng LLM}
Sau giai đoạn thử nghiệm, chúng tôi không tiếp tục sử dụng bộ phân loại zero-shot ở đầu pipeline do mô-đun truy hồi kết hợp reranker đã cho độ chính xác đủ cao để chuyển trực tiếp sang bài toán hỏi đáp. Pipeline hiện tại được tổ chức theo hướng truy hồi tăng cường ngữ cảnh cho LLM. Với mỗi câu hỏi đánh giá, hệ thống trước hết thực hiện phân rã câu hỏi và chuẩn hóa các tín hiệu pháp lý quan trọng, đặc biệt là các trường hợp sửa đổi, thay thế, bãi bỏ và mốc thời gian hiệu lực. Sau đó, hệ thống lấy ra top-$k$ văn bản liên quan, với thiết lập mặc định là $k=5$, rồi chuyển các \texttt{uid} tương ứng sang mô-đun tổng hợp ngữ cảnh để sinh câu trả lời cuối cùng bằng LLM.

Hai thành phần đánh giá được tách riêng thành pha trực tuyến và ngoại tuyến. Trong pha trực tuyến, tệp \texttt{eval\_qa\_online.py} đọc bộ dữ liệu đánh giá, kiểm tra điều kiện thu hồi tài liệu tham chiếu, sau đó truy xuất nội dung luật và thông tin bổ sung từ Neo4j để tạo payload ngoại tuyến. Cách tách pha này giúp giảm chi phí truy xuất lặp lại khi so sánh nhiều mô hình. Trong pha ngoại tuyến, tệp \texttt{eval\_qa\_offline.py} nạp lại payload đã chuẩn bị, tạo prompt hỏi đáp, gọi đồng thời nhiều mô hình qua backend \texttt{vLLM}, và ghi kết quả theo từng dòng để thuận tiện cho việc so sánh định lượng.

\subsection{Chuẩn bị dữ liệu và phạm vi đánh giá}
Bộ dữ liệu đánh giá được xây dựng từ hai nguồn chính. Thứ nhất, tập \texttt{QA\_NLP} được rà soát và gán nhãn lại theo văn bản pháp luật cập nhật, trong đó các mẫu liên quan đến xử phạt giao thông được chuyển sang bối cảnh của Nghị định 168/2024/NĐ-CP bằng hỗ trợ của LLM GLM-5 bản miễn phí. Trong quá trình này, các câu hỏi thuần factual hoặc định nghĩa được loại khỏi tập trung tâm để ưu tiên các mẫu đòi hỏi suy luận trên hiệu lực văn bản, sửa đổi hoặc thay thế nội dung. Thứ hai, tập dữ liệu do thành viên Thịnh xây dựng gồm 26 câu hỏi được kiểm tra lại và hiệu chỉnh nhãn để đảm bảo tính nhất quán giữa câu hỏi, căn cứ pháp lý và đáp án chuẩn.

Phạm vi đánh giá tập trung vào những tình huống mà hệ thống hỏi đáp pháp lý thường gặp khó khăn, bao gồm: (i) câu hỏi có tham chiếu thời điểm áp dụng luật; (ii) câu hỏi yêu cầu hợp nhất điều khoản gốc với thông tin sửa đổi; và (iii) câu hỏi có khả năng bị sai lệch nếu mô hình bỏ qua trạng thái bãi bỏ của điều luật. Thiết kế này phù hợp với mục tiêu kiểm tra năng lực sử dụng văn bản pháp lý còn hiệu lực thay vì chỉ đo khả năng sao chép nguyên văn từ ngữ cảnh truy hồi.

\subsection{Thiết kế prompt và lựa chọn mô hình}
Mô-đun sinh câu trả lời sử dụng prompt few-shot với ba ví dụ minh họa trong tệp \texttt{prompt\_qa\_fewshot.md}. Các ví dụ này mô tả trực tiếp ba kiểu tình huống quan trọng của miền pháp lý Việt Nam, gồm cập nhật mức phạt sau sửa đổi, áp dụng mốc hiệu lực cụ thể, và từ chối trả lời khi ngữ cảnh không chứa đủ căn cứ. Ngoài phần điều luật truy hồi được, prompt còn đưa vào danh sách văn bản liên quan và thông tin bổ sung về sửa đổi hoặc bãi bỏ. Một lựa chọn kỹ thuật quan trọng của hệ thống là khi truy hồi được một nút điều, khoản hoặc điểm, ngữ cảnh được mở rộng để bao gồm các nút con hoặc nút cha có liên quan, từ đó giảm nguy cơ LLM trả lời thiếu căn cứ pháp lý.

Các mô hình được xem xét cho bài toán hỏi đáp gồm Gemma 4, Qwen3.5 và Nemotron 4. Quan sát sơ bộ cho thấy các mô hình Qwen có kích thước nhỏ hơn 4B không hiểu và trả lời tiếng Việt đủ tốt cho miền pháp lý, do đó không được sử dụng trong thí nghiệm chính. Tương tự, DeepSeek cũng không được chọn vì chất lượng hiểu tiếng Việt chưa đáp ứng yêu cầu của bộ dữ liệu. Bên cạnh hướng prompting, chúng tôi duy trì nhánh thử nghiệm tinh chỉnh nhẹ bằng QLoRA trong tệp \texttt{train\_qlora.py}. Cấu hình này sử dụng lượng tử hóa 4-bit, LoRA trên các projection chính của mô hình, và huấn luyện theo định dạng hội thoại để mở ra khả năng kết hợp giữa fine-tuning và chiến lược prompting ở các bước tiếp theo.

\subsection{Thiết kế đánh giá tự động}
Đánh giá tự động của hệ thống gồm hai lớp bổ sung lẫn nhau. Lớp thứ nhất là các chỉ số đối sánh văn bản, bao gồm BLEU, ROUGE-L, METEOR và BERTScore. Trong tệp \texttt{eval\_qa\_offline.py}, các chỉ số này được tính giữa câu trả lời sinh ra và đáp án chuẩn, trong đó BERTScore sử dụng \texttt{vinai/phobert-base} để phản ánh tốt hơn tương đồng ngữ nghĩa trong tiếng Việt. Nhóm chỉ số này phù hợp để theo dõi xu hướng cải thiện tổng quát, nhưng chưa đủ để đánh giá đầy đủ tính đúng đắn pháp lý của câu trả lời.

Vì vậy, lớp đánh giá thứ hai sử dụng mô hình ngôn ngữ làm bộ chấm điểm. Prompt \texttt{prompt\_eval\_qa\_fewshot.md} định nghĩa bốn tiêu chí chấm điểm: tính chính xác về mốc thời gian, độ chính xác nội dung pháp lý, tính hữu dụng của câu trả lời, và mức độ trích dẫn căn cứ pháp lý. Thiết kế few-shot của prompt này đặc biệt nhấn mạnh hai lỗi thường gặp trong miền pháp lý giao thông: không áp dụng nội dung sửa đổi mới nhất và bỏ qua việc điều luật đã bị bãi bỏ. Trong các thử nghiệm hiện tại, cấu hình few-shot với Gemma 4 cho kết quả chấm điểm ổn định nhất, do đó được chọn làm bộ đánh giá chính cho giai đoạn draft.

\subsection{Kết quả sơ bộ và nhận xét}
Các quan sát ban đầu cho thấy việc đưa thông tin sửa đổi, thay thế và mốc hiệu lực trực tiếp vào ngữ cảnh đã cải thiện đáng kể khả năng trả lời đối với các câu hỏi cần hợp nhất văn bản. Đặc biệt, bộ dữ liệu được gán nhãn lại theo Nghị định 168/2024/NĐ-CP cho phép kiểm tra rõ hơn những trường hợp điều luật cũ không còn hiệu lực nhưng vẫn xuất hiện trong kết quả truy hồi. Đây là dạng lỗi có mức độ nghiêm trọng cao trong ứng dụng pháp lý, vì một câu trả lời có vẻ đúng về mặt cú pháp vẫn có thể sai hoàn toàn về căn cứ áp dụng.

Ở mức định tính, prompt few-shot giúp câu trả lời có cấu trúc ổn định hơn và bám sát văn bản luật hơn so với thiết lập zero-shot. Tuy nhiên, chất lượng đầu ra vẫn phụ thuộc mạnh vào khả năng hiểu tiếng Việt của mô hình nền. Quan sát này giải thích vì sao các mô hình nhỏ hơn 4B hoặc các mô hình chưa được tối ưu cho tiếng Việt không phù hợp với bài toán hiện tại. Kết quả cũng cho thấy đánh giá chỉ bằng các độ đo n-gram là chưa đủ, bởi các câu trả lời pháp lý có thể diễn đạt khác nhau nhưng vẫn đúng về nội dung, trong khi một câu trả lời sao chép gần giống đáp án chuẩn vẫn có thể sai về hiệu lực văn bản nếu bỏ qua sửa đổi.

\subsection{Hướng hoàn thiện tiếp theo}
Trong giai đoạn tiếp theo, chúng tôi dự kiến mở rộng đánh giá theo ba hướng. Thứ nhất, cần đánh giá có hệ thống tác động của chain-of-thought đối với các câu hỏi yêu cầu suy luận trên sửa đổi và bãi bỏ. Thứ hai, cần so sánh các tổ hợp giữa fine-tuning và prompting, chẳng hạn fine-tuning kết hợp chain-of-thought hoặc fine-tuning kết hợp few-shot. Thứ ba, giá trị của tham số top-$k$ cần được khảo sát thêm để xác định điểm cân bằng giữa độ đầy đủ của ngữ cảnh và nhiễu thông tin khi đưa vào LLM.

Nhìn chung, kết quả hiện tại cho thấy hướng tiếp cận retrieval-augmented legal QA kết hợp few-shot prompting và LLM-as-a-judge là khả thi cho miền pháp lý tiếng Việt. Tuy vậy, để hình thành một hệ thống đủ tin cậy cho báo cáo hoàn chỉnh, cần bổ sung bảng kết quả định lượng cuối cùng, phân tích lỗi theo nhóm hiện tượng pháp lý, và so sánh chặt chẽ hơn giữa các cấu hình mô hình nền.
